\subsection{解集的描述}
具有唯一解的线性方程组，其解集只有一个元素。
无解的线性方程组，其解集是空集。
在这两种情形下，解集都很容易描述。
只有当解集包含许多元素时，描述解集才是一个挑战。

\begin{example}
这个方程组有很多解，因为化为阶梯形后
\begin{align*}
  \begin{linsys}{3}
    2x  &   &   &+  &z  &=  &3 \\
     x  &-  &y  &-  &z  &=  &1 \\
    3x  &-  &y  &   &   &=  &4 
  \end{linsys}
  &\grstep[-(3/2)\rho_1 +\rho_3]{-(1/2)\rho_1+\rho_2}
  \begin{linsys}{3}
     2x  &   &   &+  &z      &=  &3    \\
         &   &-y &-  &(3/2)z &=  &-1/2 \\
         &   &-y &-  &(3/2)z &=  &-1/2 
   \end{linsys}                                   \\
  &\grstep{-\rho_2+\rho_3}
  \begin{linsys}{3}
     2x  &   &   &+  &z      &=  &3    \\
         &   &-y &-  &(3/2)z &=  &-1/2 \\
         &   &   &   &0      &=  &0    
   \end{linsys}
\end{align*}
并非所有变量都是首变量。
\nearbytheorem{th:GaussMethod}表明，一个$(x,y,z)$
满足第一个方程组当且仅当它满足第三个方程组。
于是我们可以将解集
$\set{(x,y,z)\suchthat\text{$2x+z=3$ 且 $x-y-z=1$ 且 $3x-y=4$}}$
描述成下面的形式。
\begin{equation*}
  \set{(x,y,z)\suchthat\text{$2x+z=3$ 且 $-y-3z/2=-1/2$}}
  \tag{$*$}
\end{equation*}
这个描述更好，因为它用两个方程代替了三个方程，
但它还不是最优的，
因为变量之间仍有一些难以理解的相互作用。

为了改进它，用不引领任何方程的变量$z$来描述
起引领作用的变量$x$和$y$。
第二个方程给出
$y=(1/2)-(3/2)z$，
第一个方程给出
$x=(3/2)-(1/2)z$。
于是我们可以把解集描述成下面这个三元组的集合。
\begin{equation*}
  \set{((3/2)-(1/2)z,\,(1/2)-(3/2)z,\,z)\suchthat z\in\Re}
  \tag{$**$}
\end{equation*}
\end{example}

与($*$)相比，
($**$)的优点
是$z$可以取任何实数。
这使得判断哪些元组在解集之中的工作
容易得多。
例如，取$z=2$就表明$(1/2,-5/2,2)$是一个解。

\begin{definition} \label{df:FreeVars}
%<*df:FreeVars>
在一个阶梯形的线性方程组中，不是首变量的变量
称为
\definend{自由变量}。\index{echelon form!free variable}\index{free variable}
%</df:FreeVars>
\end{definition}

\begin{example}   \label{ex:Parametrize2}
对一个线性方程组进行消元，结束时可能有不止一个变量是自由的。
对这个方程组用高斯消元法
\begin{align*}
   \begin{linsys}{4}
               x  &+  &y   &+  &z   &-  &w   &=  &1  \\
                  &   &y   &-  &z   &+  &w   &=  &-1 \\
              3x  &   &    &+  &6z  &-  &6w  &=  &6  \\
                  &   &-y  &+  &z   &-  &w   &=  &1  
   \end{linsys}
  &\grstep{-3\rho_1 +\rho_3}
  \begin{linsys}{4}
     x  &+  &y   &+  &z   &-  &w   &=  &1  \\
        &   &y   &-  &z   &+  &w   &=  &-1 \\
        &   &-3y &+  &3z  &-  &3w  &=  &3  \\
        &   &-y  &+  &z   &-  &w   &=  &1  
  \end{linsys}                                      \\
  &\grstep[\rho_2 +\rho_4]{3\rho_2 +\rho_3}
  \begin{linsys}{4}
     x  &+  &y   &+  &z   &-  &w   &=  &1  \\
        &   &y   &-  &z   &+  &w   &=  &-1 \\
        &   &    &   &    &   &0   &=  &0  \\
        &   &    &   &    &   &0   &=  &0  
   \end{linsys}
\end{align*}
留下\( x \)和\( y \)为首变量，而\( z \)和\( w \)都是自由变量。
为了得到我们更喜欢的描述，我们从底部往上做。
我们先把首变量$y$用
$z$和$w$表示，即$y=-1+z-w$。
再往上看最上面的方程，
代入$y$得到
$x+(-1+z-w)+z-w=1$，解出$x$得$x=2-2z+2w$。
解集
\begin{equation*}
   \set{(2-2z+2w,-1+z-w,z,w)\suchthat z,w\in\Re}
  \tag{$**$}
\end{equation*}
把首变量用自由变量表示了出来。
\end{example}

\begin{example}    \label{ex:Parametrize1}
首变量的列表可能会跳过某些列。
经过这样的消元
\begin{align*}
  \begin{linsys}{4}
              2x  &-  &2y  &   &    &   &    &=  &0  \\
                  &   &    &   &z   &+  &3w  &=  &2  \\
              3x  &-  &3y  &   &    &   &    &=  &0  \\
               x  &-  &y   &+  &2z  &+  &6w  &=  &4  
  \end{linsys}
  &\grstep[-(1/2)\rho_1+ \rho_4]{-(3/2)\rho_1 +\rho_3}
  \begin{linsys}{4}
     2x  &-  &2y  &\spaceforemptycolumn   &    &   &    &=  &0  \\
         &   &    &   &z   &+  &3w  &=  &2  \\
         &   &    &   &    &   &0   &=  &0  \\
         &   &    &   &2z  &+  &6w  &=  &4  
   \end{linsys}                                    \\
  &\grstep{-2\rho_2 +\rho_4}
  \begin{linsys}{4}
     2x  &-  &2y  &\spaceforemptycolumn   &    &   &    &=  &0  \\
         &   &    &   &z   &+  &3w  &=  &2  \\
         &   &    &   &    &   &0   &=  &0  \\
         &   &    &   &    &   &0   &=  &0  
   \end{linsys}
\end{align*}
$x$和$z$是首变量，而不是$x$和$y$。
自由变量是$y$和$w$，因此我们可以把解集描述为
$\set{ (y,y,2-3w,w)\suchthat y,w\in\Re }$。
例如，\( (1,1,2,0) \)满足该方程组\Dash 取
$y=1$和$w=0$。
四元组\( (1,0,5,4) \)不是解，
因为它的第一个分量不等于第二个分量。
\end{example}

%<*df:Parameter>
我们用来描述一族解的变量是一个\definend{参数}。\index{parameter}  
%</df:Parameter>
我们说前面的例题中的解集
是用$y$和$w$
来\definend{参数化\/}\index{parametrized} 
的。

“参数”和“自由变量”这两个术语的含义并不相同。
在前面的例题中
$y$和$w$是自由变量，因为在阶梯形方程组中它们
不引领任何方程。
它们之所以是参数，是因为
我们用它们来描述解的集合。
假如我们改而
把第二个方程改写成$w=2/3-(1/3)z$，那么
自由变量仍然是$y$和$w$，但参数
将是$y$和$z$。

在本书余下的部分中
我们将通过把线性方程组化为
阶梯形、然后用自由变量进行参数化来求解线性方程组。

\begin{example}
这是另一个有无穷多解的方程组。
\begin{align*}
   \begin{linsys}{4}
               x  &+  &2y  &   &   &   &   &=  &1  \\
              2x  &   &    &+  &z  &   &   &=  &2  \\
              3x  &+  &2y  &+  &z  &-  &w  &=  &4  
   \end{linsys}
  &\grstep[-3\rho_1 +\rho_3]{-2\rho_1+\rho_2}
  \begin{linsys}{4}
     x  &+  &2y  &   &   &   &   &=  &1  \\
        &   &-4y &+  &z  &   &   &=  &0  \\
        &   &-4y &+  &z  &-  &w  &=  &1  
   \end{linsys}                                    \\
  &\grstep{-\rho_2+\rho_3}
  \begin{linsys}{4}
     x  &+  &2y  &   &   &   &   &=  &1  \\
        &   &-4y &+  &z  &   &   &=  &0  \\
        &   &    &   &   &   &-w &=  &1  
   \end{linsys}
\end{align*}
首变量是\( x \)、\( y \)和\( w \)。
变量\( z \)是自由的。
注意，虽然有无穷多
解，但$w$的值并不变化，而是恒为$w=-1$。
为了参数化，把\( w \)用\( z \)写成\( w=-1+0z \)。
于是\( y=(1/4)z \)。
把\( y \)代入第一个
方程得到\( x=1-(1/2)z \)。
解集是$\set{(1-(1/2)z,(1/4)z,z,-1)\suchthat z\in\Re}$。
\end{example}

对解集进行参数化表明，带自由变量的方程组
有无穷多解。
例如，上面的$z$取遍无穷多个实数值，
每个值对应一个不同的解。

我们以建立一套简化的记法来结束本小节，
这套记法用于线性方程组
及其解集。

\begin{definition}\label{df:matrix}
%<*df:matrix>
一个\( \nbym{m}{n} \) \definend{矩阵}\index{matrix}
是数的矩形阵列，
有\( m \)~\definend{行}\index{matrix!row}\index{row} 
和\( n \)~\definend{列}\index{matrix!column}\index{column}。
矩阵中的每个数都是一个
\definend{元素}\index{matrix!entry}\index{entry, matrix}。
%</df:matrix>
\end{definition}

我们通常用一个大写罗马字母来表示矩阵。
例如，
\begin{equation*}
  A=
  \begin{mat}[r]
    1  &2.2  &5  \\
    3  &4    &-7
  \end{mat}
\end{equation*}
它有$2$~行和$3$~列，因此
是一个\( \nbym{2}{3} \)矩阵。
把它读作“二乘三”；
行的个数总是先说。
（矩阵围有括号，
这样当
两个矩阵相邻时，
我们能分辨出一个在哪里结束、另一个从哪里开始。）
我们用对应的小写字母命名矩阵的元素，
于是上面阵列中第二行第一列的
元素是\( a_{2,1}=3 \)。
注意下标的次序很重要：
$a_{1,2}\neq a_{2,1}$，因为\( a_{1,2}=2.2 \)。
我们把
所有\( \nbym{m}{n} \)
矩阵组成的集合记作\( \matspace_{\nbym{m}{n}} \)。\index{matrices|set of}

我们用矩阵做高斯消元法，其方式与
我们对方程组所用的方式基本相同：
矩阵一行的
\definend{首主元}\index{echelon form!leading entry}%
\index{leading!entry} %
是它的第一个非零元素（如果有的话），我们进行行变换来达到
\definend{矩阵阶梯形}，\index{echelon form!matrix}%
\index{matrix!echelon form} %
即下面各行的首主元位于其上方各
行首主元的右侧。
我们喜欢矩阵记法，因为它减轻了
抄写工作的负担，即抄写变量以及
书写$+$号和$=$号的负担。

\begin{example}
我们可以将这个线性方程组
\begin{equation*}
  \begin{linsys}{3}
    x  &+  &2y  &   &    &=  &4   \\
       &   &y   &-  &z   &=  &0   \\
    x  &   &    &+  &2z  &=  &4   
  \end{linsys}
\end{equation*}
缩写成这个矩阵。
\begin{equation*}
    \begin{amat}[r]{3}
      1  &2  &0  &4  \\
      0  &1  &-1 &0  \\
      1  &0  &2  &4
    \end{amat}
\end{equation*}
竖线提醒读者方程组左边的
系数与右边的常数之间的区别。
带竖线的矩阵是
\definend{增广\/}\index{matrix!augmented}\index{augmented matrix}矩阵。
\begin{equation*}
    \begin{amat}[r]{3}
      1  &2  &0  &4  \\
      0  &1  &-1 &0  \\
      1  &0  &2  &4
    \end{amat}
  \grstep{-\rho_1 +\rho_3}
  \begin{amat}[r]{3}
       1  &2  &0  &4  \\
       0  &1  &-1 &0  \\
       0  &-2 &2  &0
     \end{amat}                        
  \grstep{2\rho_2 +\rho_3}
  \begin{amat}[r]{3}
       1  &2  &0  &4  \\
       0  &1  &-1 &0  \\
       0  &0  &0  &0
     \end{amat}
\end{equation*}
第二行表示$y-z=0$，第一行表示
$x+2y=4$，因此解集是
\( \set{(4-2z,z,z)\suchthat z\in\Re} \)。
\end{example}

矩阵记法也
使解集的描述更加清晰。
\nearbyexample{ex:Parametrize2}的
$\set{(2-2z+2w,-1+z-w,z,w)\suchthat z,w\in\Re}$ 
很难读。
我们将把它改写，使所有常数放在一起，\( z \)的所有系数放在一起，以及\( w \)的所有系数
放在一起。
我们把它们竖着写成一列的矩阵。
\begin{equation*}
  \set{\colvec[r]{2 \\ -1 \\ 0 \\ 0}
       +\colvec[r]{-2 \\ 1 \\ 1 \\ 0}\cdot z
       +\colvec[r]{2 \\ -1 \\ 0 \\ 1}\cdot w
       \suchthat z,w\in\Re}
\end{equation*}
例如，最上面一行说明\( x=2-2z+2w \)，
第二行说明\( y= -1+z-w \)。
（下一节给出一个几何解释，它将帮助
我们想象解集。）

\begin{definition}\label{df:vector}
%<*df:vector>
一个
\definend{列向量}，\index{column!vector}\index{vector!column}
通常简称\definend{向量}，\index{vector} 
是只有一列的矩阵。
只有一行的矩阵是
\definend{行向量}\index{row!vector}\index{vector!row}。
向量的元素有时称为
\definend{分量}\index{component of a vector}\index{vector!component}。
分量全为零的列向量或行向量是一个
\definend{零向量}。\index{zero vector}\index{vector!zero}
%</df:vector>
\end{definition}

向量是用大写罗马字母表示矩阵这一惯例的一个例外。
我们使用带上箭头的小写罗马字母或希腊字母：
\( \vec{a} \)、\( \vec{b} \)、\ldots\, 或
\( \vec{\alpha} \)、\( \vec{\beta} \)、\ldots\
（粗体也很常见：
{\boldmath \( a \)}或{\boldmath \( \alpha \)}）。
例如，这是一个第三个分量为\( 7 \)的列向量。
\begin{equation*}
  \vec{v}=
  \colvec[r]{ 1  \\  3  \\ 7}
\end{equation*}
零向量记作\( \zero \)。
零向量有很多种\Dash 高度为一的零向量、高度为二的零向量，等等\Dash 但
尽管如此，我们仍常说“这个”零向量，
并期望由上下文可以清楚地知道其大小。

\begin{definition}
带有未知数\( x_1,\ldots\,,x_n \)的线性方程
\( a_1x_1+\,\cdots\,+a_nx_n=d \)
被
\begin{equation*}
  \vec{s}=\colvec{s_1 \\ \vdotswithin{s_1} \\ s_n}
\end{equation*}
\definend{满足}\index{vector!satisfies an equation}%
\index{linear equation!satisfied by a vector}，是指
\( a_1s_1+\,\cdots\,+a_ns_n=d \)。
若一个向量满足线性方程组中的每一个方程，则称它满足该线性方程组。
\end{definition}

我们所使用的解集描述方式
涉及把向量相加，
以及用实数去乘它们。
在给出展示这种描述方式的例子之前
我们首先需要定义这两种运算。

\begin{definition} \label{df:VectorSum}
%<*df:VectorSum>
\( \vec{u} \)与\( \vec{v} \)的\definend{向量和}\index{vector!sum}\index{sum!vector}\index{addition of vectors} 
是由各分量之和组成的向量。
\begin{equation*}
  \vec{u}+\vec{v}=
  \colvec{u_1 \\ \vdotswithin{u_1} \\ u_n}
   +
  \colvec{v_1 \\ \vdotswithin{v_1} \\ v_n}
   =
  \colvec{u_1+v_1 \\ \vdotswithin{u_1+v_1} \\ u_n+v_n}
\end{equation*}
%</df:VectorSum>
\end{definition}

注意，要使加法有定义，
各向量必须有相同个数的元素。
这种逐元素的加法对任意一对矩阵都适用，而不只是对向量，
只要它们有相同的行数和
列数。\index{matrix!sum}\index{sum!matrix}

\begin{definition} \label{df:VectorScalarMultiplication}
%<*df:VectorScalarMultiplication>
实数
\( r \)与向量\( \vec{v} \)的\definend{标量乘法\/}\index{vector!scalar multiple}%
\index{scalar multiple!vector}是由各倍数组成的向量。
\begin{equation*}
  r\cdot\vec{v}=
  r\cdot\colvec{v_1 \\ \vdotswithin{v_1} \\ v_n}
  =
  \colvec{rv_1 \\ \vdotswithin{rv_1} \\ rv_n}
\end{equation*}
%</df:VectorScalarMultiplication>
\end{definition}

与加法运算一样，逐元素的标量乘法
运算也超出向量的范围而适用于任意
矩阵。\index{matrix!scalar multiplication}\index{scalar multiplication!matrix}

标量乘法可以写成\( r\cdot\vec{v} \)或
\( \vec{v}\cdot r \)，有时甚至省略“$\cdot$”符号：~$r\vec{v}$。
（不要把标量乘法
称为“标量积”，因为那个名称指的是另一种运算。）

\begin{example}
\begin{equation*}
  \colvec[r]{2 \\ 3 \\ 1}
   +
  \colvec[r]{3 \\ -1 \\ 4}
  =
  \colvec{2+3 \\ 3-1 \\ 1+4}
   =
  \colvec[r]{5 \\ 2 \\ 5}
  \qquad
  7\cdot\colvec[r]{1 \\ 4 \\ -1 \\ -3}
  =
  \colvec[r]{7 \\ 28 \\ -7 \\ -21}
\end{equation*}
\end{example}

注意，加法与标量乘法的定义在重合之处是一致的；例如，\( \vec{v} +\vec{v} = 2\vec{v} \)。

有了这些定义，我们就可以用矩阵与向量记法
既求解方程组又表示解了。

\begin{example} \label{ex:ManyParamsInfManySolsSystem}
这个方程组
\begin{equation*}
   \begin{linsys}{5}
      2x  &+  &y  &  &  &-  &w  &   &   &=  &4  \\
          &   &y  &  &  &+  &w  &+  &u  &=  &4  \\
       x  &   &   &- &z &+  &2w &   &   &=  &0  
   \end{linsys}
\end{equation*}
按如下方式化简。
\begin{align*}
  \begin{amat}[r]{5}
    2  &1  &0  &-1  &0  &4  \\
    0  &1  &0  &1   &1  &4  \\
    1  &0  &-1 &2   &0  &0
  \end{amat}
  &\grstep{-(1/2)\rho_1+\rho_3}
  \begin{amat}[r]{5}
    2  &1     &0  &-1    &0  &4  \\
    0  &1     &0  &1     &1  &4  \\
    0  &-1/2  &-1 &5/2   &0  &-2
  \end{amat}                                 \\
  &\grstep{(1/2)\rho_2+\rho_3}
  \begin{amat}[r]{5}
    2  &1     &0  &-1    &0    &4  \\
    0  &1     &0  &1     &1    &4  \\
    0  &0     &-1 &3     &1/2  &0
  \end{amat}
\end{align*}
解集是
\( \set{(w+(1/2)u,4-w-u,3w+(1/2)u,w,u)\suchthat w,u\in\Re} \)。
我们把它写成向量形式。
\begin{equation*}
  \set{\colvec{x \\ y \\ z \\ w \\ u}=
       \colvec[r]{0 \\ 4 \\ 0 \\ 0 \\ 0}+
       \colvec[r]{1 \\ -1 \\ 3 \\ 1 \\ 0}w+
       \colvec[r]{1/2 \\ -1 \\ 1/2 \\ 0 \\ 1}u
       \suchthat w,u\in\Re}
\end{equation*}
注意向量记法把
每个参数的系数区分得多么清楚。
例如，向量形式的第三行清楚地表明，如果\( u \)固定，那么\( z \)增长的速度是\( w \)的三倍。
同样清楚表明的是，令\( w \)和\( u \)都为零
得到
\begin{equation*}
  \colvec{x \\ y \\ z \\ w \\ u}
  =\colvec[r]{0 \\ 4 \\ 0 \\ 0 \\ 0}
\end{equation*}
是线性方程组的一个特解。
\end{example}

\begin{example}
同样地，方程组
\begin{equation*}
   \begin{linsys}{3}
     x  &-  &y  &+  &z  &=  &1  \\
    3x  &   &   &+  &z  &=  &3  \\
    5x  &-  &2y &+  &3z &=  &5  
  \end{linsys}
\end{equation*}
化简
\begin{align*}
  \begin{amat}[r]{3}
    1  &-1  &1  &1  \\
    3  &0   &1  &3  \\
    5  &-2  &3  &5
  \end{amat}
  &\grstep[-5\rho_1+\rho_3]{-3\rho_1+\rho_2}
  \begin{amat}[r]{3}
    1  &-1  &1  &1  \\
    0  &3   &-2 &0  \\
    0  &3   &-2 &0
  \end{amat}                                    \\
  &\grstep{-\rho_2+\rho_3}
  \begin{amat}[r]{3}
    1  &-1  &1  &1  \\
    0  &3   &-2 &0  \\
    0  &0   &0  &0
  \end{amat}
\end{align*}
后给出一个单参数的解集。
\begin{equation*}
  \set{\colvec[r]{1 \\ 0 \\ 0}
       +\colvec[r]{-1/3 \\ 2/3 \\ 1}z
       \suchthat z\in\Re}
\end{equation*}
与前面的例题中一样，与参数无关的向量
\begin{equation*}
   \colvec[r]{1 \\ 0 \\ 0}
\end{equation*}
是该方程组的一个特解。
\end{example}

在练习之前，我们先考虑一下已经完成了什么，
以及本章余下部分要做什么。
到目前为止，我们做的是高斯消元法的机械步骤。
我们还没有停下来考虑过出现的任何问题，
只是证明了\nearbytheorem{th:GaussMethod}\Dash 该定理
通过表明这个方法给出
正确的答案而为方法提供了依据。

例如，我们总能像上面那样描述解集吗，
即用一个特解向量
加上若干其他向量的不受限制的线性组合？
我们已经注意到，这样描述的解集
有无穷多个成员，
因此回答这个问题
会告诉我们解集的大小。
下一小节将表明答案是“肯定的”。
本章第二节随后
将利用这个答案来描述解集的几何。

另一些问题来自这样的观察：我们做高斯消元法可以
不止一种方式（例如，对换行时，我们可能有多行
可供选择对换）。
\nearbytheorem{th:GaussMethod}说，无论我们怎么做，都必定得到同一个解集，但是
如果用两种方式做高斯消元法，
在每个阶梯形方程组中我们必定得到相同个数的自由变量吗？
那些必定是同样的变量吗？也就是说，用一种方式求解一个问题得到$y$和$w$为自由变量，而用另一种方式求解却得到$y$和$z$为自由变量，这是不可能的吗？
本章第三节回答“是的”：
从任何初始线性方程组出发，
所有导出的阶梯形版本
都有相同的自由变量。

因此，到本章结束时，我们不仅会在
高斯消元法的实践上有扎实的基础，
也会在理论上
有扎实的基础。
我们将确切地知道在一次消元中什么会发生、什么不会发生。

\begin{exercises}
  \recommended \item  
    求矩阵的指定元素，
    （若它有定义）。
    \begin{equation*}
      A=\begin{mat}[r]
        1  &3  &1  \\
        2  &-1 &4
      \end{mat}
    \end{equation*}
    \begin{exparts*}
      \partsitem \( a_{2,1} \)
      \partsitem \( a_{1,2} \)
      \partsitem \( a_{2,2} \)
      \partsitem \( a_{3,1} \)
    \end{exparts*}
    \begin{answer}
      \begin{exparts*}
        \partsitem \( 2 \)
        \partsitem \( 3 \)
        \partsitem \(-1 \)
        \partsitem 无定义。
      \end{exparts*}  
    \end{answer}
  \recommended \item 
    给出每个矩阵的尺寸。
    \begin{exparts*}
      \partsitem \(
        \begin{mat}[r]
          1  &0  &4  \\
          2  &1  &5
        \end{mat}  \)
      \partsitem \(
        \begin{mat}[r]
          1  &1  \\
         -1  &1  \\
          3  &-1
        \end{mat}  \)
      \partsitem \(
        \begin{mat}[r]
          5  &10 \\
         10  &5
        \end{mat}  \)
    \end{exparts*}
    \begin{answer}
      \begin{exparts*}
        \partsitem \( \nbym{2}{3} \)
        \partsitem \( \nbym{3}{2} \)
        \partsitem \( \nbym{2}{2} \)
      \end{exparts*}  
    \end{answer}
  \recommended \item 
    进行指定的向量运算（若它有定义）。
    \begin{exparts*}
      \partsitem \( \colvec[r]{2 \\ 1 \\ 1}
               +\colvec[r]{3 \\ 0 \\ 4} \)
      \partsitem \( 5\colvec[r]{4 \\ -1} \)
      \partsitem \( \colvec[r]{1 \\ 5 \\ 1}
               -\colvec[r]{3 \\ 1 \\ 1} \)
      \partsitem \( 7\colvec[r]{2 \\ 1}
               +9\colvec[r]{3 \\ 5} \)
      \partsitem \( \colvec[r]{1 \\ 2}
               +\colvec[r]{1 \\ 2 \\ 3} \)
      \partsitem \( 6\colvec[r]{3 \\ 1 \\ 1}
               -4\colvec[r]{2 \\ 0 \\ 3}
               +2\colvec[r]{1 \\ 1 \\ 5} \)
    \end{exparts*}
    \begin{answer}
      \begin{exparts*}
        \partsitem \( \colvec[r]{5 \\ 1 \\ 5} \)
        \partsitem \( \colvec[r]{20 \\ -5} \)
        \partsitem \( \colvec[r]{-2 \\ 4 \\ 0} \)
        \partsitem \( \colvec[r]{41 \\ 52} \)
        \partsitem 无定义。
        \partsitem \( \colvec[r]{12 \\ 8 \\ 4} \)
      \end{exparts*}  
     \end{answer}
  \recommended \item  \label{exer:SolveInMatrixNotation}
    用矩阵记法求解每个方程组。
    用向量表示解。
    \begin{exparts*}
      \partsitem \( \begin{linsys}[t]{2}
                  3x  &+  &6y  &=  &18  \\
                   x  &+  &2y  &=  &6   
                   \end{linsys}  \)
      \partsitem \( \begin{linsys}[t]{2}
                   x  &+  &y   &=  &1  \\
                   x  &-  &y   &=  &-1   
                    \end{linsys}  \)
      \partsitem \( \begin{linsys}[t]{3}
                   x_1  &   &     &+  &x_3   &=  &4  \\
                   x_1  &-  &x_2  &+  &2x_3  &=  &5  \\
                  4x_1  &-  &x_2  &+  &5x_3  &=  &17  
                   \end{linsys}  \)
      \partsitem \( \begin{linsys}[t]{3}
                   2a   &+  &b    &-  &c     &=  &2  \\
                   2a   &   &     &+  &c     &=  &3  \\
                    a   &-  &b    &   &      &=  &0   
                    \end{linsys}  \)
      \partsitem \( \begin{linsys}[t]{4}
                     x  &+  &2y   &-   &z   &    &    &=  &3  \\
                    2x  &+  &y    &    &    &+   &w   &=  &4  \\
                     x  &-  &y    &+   &z   &+   &w   &=  &1  
                    \end{linsys}  \)
      \partsitem \( \begin{linsys}[t]{4}
                     x  &   &     &+   &z   &+   &w   &=  &4  \\
                    2x  &+  &y    &    &    &-   &w   &=  &2  \\
                    3x  &+  &y    &+   &z   &    &    &=  &7  
                     \end{linsys}  \)
    \end{exparts*}
    \begin{answer}
      \begin{exparts}
        \partsitem 这一消元
          \begin{equation*}
            \begin{amat}[r]{2}
              3  &6  &18 \\
              1  &2  &6
            \end{amat}
            \grstep{(-1/3)\rho_1+\rho_2}
            \begin{amat}[r]{2}
              3  &6  &18 \\
              0  &0  &0
            \end{amat}
          \end{equation*}
          留下\( x \)为首变量而\( y \)为自由变量。
          取\( y \)为参数，得到\( x=6-2y \)以及这个解
          集。
          \begin{equation*}
            \set{\colvec[r]{6 \\ 0}+\colvec[r]{-2 \\ 1}y
              \suchthat y\in\Re}
          \end{equation*}
        \partsitem 一次消元
          \begin{equation*}
            \begin{amat}[r]{2}
              1  &1  &1  \\
              1  &-1 &-1
            \end{amat}
            \grstep{-\rho_1+\rho_2}
            \begin{amat}[r]{2}
              1  &1  &1  \\
              0  &-2 &-2
            \end{amat}
          \end{equation*}
          给出唯一解\( y=1 \)、\( x=0 \)。
          解集如下。
          \begin{equation*}
            \set{\colvec[r]{0 \\ 1} }
          \end{equation*}
        \partsitem 高斯消元法
          \begin{equation*}
            \begin{amat}[r]{3}
              1  &0  &1  &4  \\
              1  &-1 &2  &5  \\
              4  &-1 &5  &17
            \end{amat}
            \grstep[-4\rho_1+\rho_3]{-\rho_1+\rho_2}
            \begin{amat}[r]{3}
              1  &0  &1  &4  \\
              0  &-1 &1  &1  \\
              0  &-1 &1  &1
            \end{amat}     
            \grstep{-\rho_2+\rho_3}
            \begin{amat}[r]{3}
              1  &0  &1  &4  \\
              0  &-1 &1  &1  \\
              0  &0  &0  &0
            \end{amat}
          \end{equation*}
          留下\( x_1 \)和\( x_2 \)为首变量，而\( x_3 \)为自由变量。
          解集如下。
          \begin{equation*}
            \set{\colvec[r]{4 \\ -1 \\ 0}+\colvec[r]{-1 \\ 1 \\ 1}x_3
              \suchthat x_3\in\Re}
          \end{equation*}
        \partsitem 这一消元
          \begin{align*}
            \begin{amat}[r]{3}
              2  &1  &-1 &2  \\
              2  &0  &1  &3  \\
              1  &-1 &0  &0
            \end{amat}
            &\grstep[-(1/2)\rho_1+\rho_3]{-\rho_1+\rho_2}
            \begin{amat}[r]{3}
              2  &1    &-1   &2  \\
              0  &-1   &2    &1  \\
              0  &-3/2 &1/2  &-1
            \end{amat}                                          \\
            &\grstep{(-3/2)\rho_2+\rho_3}
            \begin{amat}[r]{3}
              2  &1  &-1   &2  \\
              0  &-1 &2    &1  \\
              0  &0  &-5/2 &-5/2
            \end{amat}
          \end{align*}
          表明解集是单元素集。
          \begin{equation*}
            \set{\colvec[r]{1 \\ 1 \\ 1}}
          \end{equation*}
        \partsitem 这一消元很容易
          \begin{align*}
            \begin{amat}[r]{4}
              1  &2  &-1 &0  &3 \\
              2  &1  &0  &1  &4 \\
              1  &-1 &1  &1  &1
            \end{amat}
            &\grstep[-\rho_1+\rho_3]{-2\rho_1+\rho_2}
            \begin{amat}[r]{4}
              1  &2  &-1 &0  &3  \\
              0  &-3 &2  &1  &-2 \\
              0  &-3 &2  &1  &-2
            \end{amat}                                       \\
            &\grstep{-\rho_2+\rho_3}
            \begin{amat}[r]{4}
              1  &2  &-1 &0  &3  \\
              0  &-3 &2  &1  &-2 \\
              0  &0  &0  &0  &0
            \end{amat}
          \end{align*}
          结束时\( x \)和$y$为首变量，而\( z \)和\( w \)为
          自由变量。
          解出\( y \)得到\( y=(2+2z+w)/3 \)，代入表明
          \( x+2(2+2z+w)/3-z=3 \)，于是\( x=(5/3)-(1/3)z-(2/3)w \)，
          由此得到这个解集。
          \begin{equation*}
            \set{\colvec[r]{5/3 \\ 2/3 \\ 0 \\ 0}
                 +\colvec[r]{-1/3 \\ 2/3 \\ 1 \\ 0}z
                 +\colvec[r]{-2/3 \\ 1/3 \\ 0 \\ 1}w
                 \suchthat z,w\in\Re}
          \end{equation*}
        \partsitem 消元
          \begin{align*}
            \begin{amat}[r]{4}
              1  &0  &1  &1  &4 \\
              2  &1  &0  &-1 &2 \\
              3  &1  &1  &0  &7
            \end{amat}
            &\grstep[-3\rho_1+\rho_3]{-2\rho_1+\rho_2}
            \begin{amat}[r]{4}
              1  &0  &1  &1  &4 \\
              0  &1  &-2 &-3 &-6\\
              0  &1  &-2 &-3 &-5
            \end{amat}                                       \\
            &\grstep{-\rho_2+\rho_3}
            \begin{amat}[r]{4}
              1  &0  &1  &1  &4 \\
              0  &1  &-2 &-3 &-6\\
              0  &0  &0  &0  &1
            \end{amat}
          \end{align*}
          表明无解\Dash 解集为空。
      \end{exparts}  
     \end{answer}
  \item \label{exer:SlvMatNot}
    用矩阵记法求解每个方程组。
    用向量记法给出每个解集。
    \begin{exparts*}
      \partsitem \( \begin{linsys}[t]{3}
                  2x  &+  &y  &-  &z  &=  &1  \\
                  4x  &-  &y  &   &   &=  &3  
                \end{linsys}  \)
      \partsitem \( \begin{linsys}[t]{4}
                   x  &   &   &-  &z  &   &   &=  &1  \\
                      &   &y  &+  &2z &-  &w  &=  &3  \\
                   x  &+  &2y &+  &3z &-  &w  &=  &7  
               \end{linsys}  \)
      \partsitem \( \begin{linsys}[t]{4}
                   x  &-  &y  &+  &z  &   &   &=  &0  \\
                      &   &y  &   &   &+  &w  &=  &0  \\
                  3x  &-  &2y &+  &3z &+  &w  &=  &0  \\
                      &   &-y &   &   &-  &w  &=  &0  
               \end{linsys}  \)
      \partsitem \( \begin{linsys}[t]{5}
                   a  &+  &2b &+  &3c &+  &d  &-  &e  &=  &1  \\
                  3a  &-  &b  &+  &c  &+  &d  &+  &e  &=  &3  
               \end{linsys}  \)
    \end{exparts*}
    \begin{answer}
      \begin{exparts}
      \partsitem 这一消元
        \begin{equation*}
          \begin{amat}[r]{3}
            2  &1  &-1  &1  \\
            4  &-1 &0   &3
          \end{amat}
          \grstep{-2\rho_1+\rho_2}
          \begin{amat}[r]{3}
            2  &1  &-1  &1  \\
            0  &-3 &2   &1
          \end{amat}
        \end{equation*}
        结束时\( x \)和\( y \)为首变量，而\( z \)为自由变量。
        解出\( y \)得到\( y=(1-2z)/(-3) \)，然后代入
        \( 2x+(1-2z)/(-3)-z=1 \)表明\( x=((4/3)+(1/3)z)/2 \)。
        因此解集如下。
        \begin{equation*}
          \set{\colvec[r]{2/3 \\ -1/3 \\ 0}
               +\colvec[r]{1/6 \\ 2/3 \\ 1}z
              \suchthat z\in\Re}
        \end{equation*}
      \partsitem 这样应用高斯消元法
        \begin{align*}
          \begin{amat}[r]{4}
            1  &0  &-1  &0  &1 \\
            0  &1  &2   &-1 &3 \\
            1  &2  &3   &-1 &7
          \end{amat}
          &\grstep{-\rho_1+\rho_3}
          \begin{amat}[r]{4}
            1  &0  &-1  &0  &1 \\
            0  &1  &2   &-1 &3 \\
            0  &2  &4   &-1 &6
          \end{amat}                                   \\
          &\grstep{-2\rho_2+\rho_3}
          \begin{amat}[r]{4}
            1  &0  &-1  &0  &1 \\
            0  &1  &2   &-1 &3 \\
            0  &0  &0   &1  &0
          \end{amat}
        \end{align*}
        留下\( x \)、\( y \)和\( w \)为首变量。
        解集如下。
        \begin{equation*}
          \set{\colvec[r]{1 \\ 3 \\ 0 \\ 0}
               +\colvec[r]{1 \\ -2 \\ 1 \\ 0}z
              \suchthat z\in\Re}
        \end{equation*}
      \partsitem 这次行化简
        \begin{align*}
          \begin{amat}[r]{4}
            1  &-1 &1   &0  &0 \\
            0  &1  &0   &1  &0 \\
            3  &-2 &3   &1  &0 \\
            0  &-1 &0   &-1 &0
          \end{amat}
          &\grstep{-3\rho_1+\rho_3}
          \begin{amat}[r]{4}
            1  &-1 &1   &0  &0 \\
            0  &1  &0   &1  &0 \\
            0  &1  &0   &1  &0 \\
            0  &-1 &0   &-1 &0
          \end{amat}                                       \\
          &\grstep[\rho_2+\rho_4]{-\rho_2+\rho_3}
          \begin{amat}[r]{4}
            1  &-1 &1   &0  &0 \\
            0  &1  &0   &1  &0 \\
            0  &0  &0   &0  &0 \\
            0  &0  &0   &0  &0
          \end{amat}
        \end{align*}
        结束时\( z \)和\( w \)为自由变量。
        我们得到这个解集。
        \begin{equation*}
          \set{\colvec[r]{0 \\ 0 \\ 0 \\ 0}
               +\colvec[r]{-1 \\ 0 \\ 1 \\ 0}z
               +\colvec[r]{-1 \\ -1 \\ 0 \\ 1}w
              \suchthat z,w\in\Re}
        \end{equation*}
      \partsitem 按这种方式做高斯消元法
        \begin{equation*}
          \begin{amat}[r]{5}
            1  &2  &3   &1  &-1 &1  \\
            3  &-1 &1   &1  &1  &3
          \end{amat}
          \grstep{-3\rho_1+\rho_2}
          \begin{amat}[r]{5}
            1  &2  &3   &1  &-1 &1  \\
            0  &-7 &-8  &-2 &4  &0
          \end{amat}
        \end{equation*}
        结束时\( c \)、\( d \)和\( e \)为自由变量。
        解出\( b \)表明\( b=(8c+2d-4e)/(-7) \)，然后 
        代入
        \( a+2(8c+2d-4e)/(-7)+3c+1d-1e=1 \)表明
        \( a=1-(5/7)c-(3/7)d-(1/7)e \)，于是我们得到解集。
        \begin{equation*}
          \set{\colvec[r]{1 \\ 0 \\ 0 \\ 0 \\ 0}
               +\colvec[r]{-5/7 \\ -8/7 \\ 1 \\ 0 \\ 0}c
               +\colvec[r]{-3/7 \\ -2/7 \\ 0 \\ 1 \\ 0}d
               +\colvec[r]{-1/7 \\ 4/7 \\ 0 \\ 0 \\ 1}e
              \suchthat c,d,e\in\Re}
        \end{equation*}
    \end{exparts}  
   \end{answer}
  \item 
    用矩阵记法求解每个方程组。
    用向量表示解集。
    \begin{exparts*}
      \partsitem 
        $\begin{linsys}{3}
          3x &+ &2y &+ &z &= &1 \\
          x  &- &y  &+ &z &= &2 \\
          5x &+ &5y &+ &z &= &0
        \end{linsys}$
      \partsitem
        $\begin{linsys}{4}
          x  &+ &y  &- &2z &= &0 \\
          x  &- &y  &  &   &= &-3 \\
          3x &- &y  &- &2z &= &-6  \\
             &  &2y &- &2z &= &3  
        \end{linsys}$
      \partsitem
        $\begin{linsys}{5}
          2x  &- &y  &- &z &+ &w &= &4 \\
           x  &+ &y  &+ &z &  &  &= &-1 
        \end{linsys}$
      \partsitem
         $\begin{linsys}{3}
          x  &+ &y  &- &2z &= &0 \\
          x  &- &y  &  &   &= &-3 \\
          3x &- &y  &- &2z &= &0    
        \end{linsys}$ 
    \end{exparts*}
    \begin{answer}
      \begin{exparts}
        \partsitem 这一消元
          \begin{align*}
            \begin{amat}{3}
              3 &2  &1  &1 \\
              1 &-1 &1  &2 \\
              5 &5  &1  &0
            \end{amat}
            &\grstep[-(5/3)\rho_1+\rho_3]{-(1/3)\rho_1+\rho_2}
            \begin{amat}{3}
              3 &2    &1     &1 \\
              0 &-5/3 &2/3   &5/3 \\
              0 &5/3  &-2/3  &-5/3
            \end{amat}                                     \\
            &\grstep{\rho_2+\rho_3}
            \begin{amat}{3}
              3 &2    &1     &1 \\
              0 &-5/3 &2/3   &5/3 \\
              0 &0    &0     &0
            \end{amat}
          \end{align*}
          给出这个解集。
          \begin{equation*}
            \set{\colvec{x \\ y \\ z}=\colvec{1 \\ -1 \\ 0}
                                       +\colvec{-3/5 \\ 2/5 \\ 1}z
                             \suchthat z\in\Re}
          \end{equation*}
        \partsitem
          消元如下。
          \begin{align*}
            \begin{amat}{3}
              1 &1   &-2 &0 \\
              1 &-1  &0  &3  \\
              3 &-1  &-2 &-6 \\
              0 &2   &-2 &3
            \end{amat}
            &\grstep[-3\rho_1+\rho_3]{-\rho_1+\rho_2}
            \begin{amat}{3}
              1 &1   &-2 &0 \\
              0 &-2  &2  &-3  \\
              0 &-4  &4  &-6 \\
              0 &2   &-2 &3
            \end{amat}                                  \\
            &\grstep[\rho_2+\rho_4]{-2\rho_2+\rho_3}
            \begin{amat}{3}
              1 &1   &-2 &0 \\
              0 &-2  &2  &-3  \\
              0 &0   &0  &0 \\
              0 &0   &0  &0
            \end{amat}
          \end{align*}
          解集如下。
          \begin{equation*}
            \set{\colvec{-3/2 \\ 3/2 \\ 0}
                 +\colvec{1 \\ 1 \\ 1}z
                 \suchthat z\in\Re}
          \end{equation*}
      \partsitem
         高斯消元法
         \begin{equation*}
           \begin{amat}{4}
             2 &-1 &-1 &1 &4 \\
             1 &1  &1  &0 &-1
           \end{amat}
           \grstep{-(1/2)\rho_1+\rho_2}
           \begin{amat}{4}
             2 &-1   &-1   &1    &4 \\
             0 &3/2  &3/2  &-1/2 &-3
           \end{amat}
         \end{equation*}
         给出解集。
         \begin{equation*}
           \set{\colvec{1 \\ -2 \\ 0 \\ 0}
                  +\colvec{0 \\ -1 \\ 1 \\ 0}z
                   \colvec{-1/3 \\ 1/3 \\ 0 \\ 1}w
                 \suchthat z,w\in\Re}
         \end{equation*}
       \partsitem
          消元如下。
          \begin{equation*}
            \begin{amat}{3}
              1 &1  &-2 &0  \\
              1 &-1 &0  &-3 \\
              3 &-1 &-2 &0
           \end{amat}
           \grstep[-3\rho_1+\rho_3]{-\rho_1+\rho_2}
           \begin{amat}{3}
             1 &1  &-2 &0  \\
             0 &-2 &2  &-3 \\
             0 &-4  &4  &0
          \end{amat}
          \grstep{-2\rho_2+\rho_3}
          \begin{amat}{3}
            1 &1  &-2 &0  \\
            0 &-2 &2  &-3 \\
            0 &0  &0  &6
          \end{amat}
        \end{equation*}
        解集为空集~$\set{}$。
      \end{exparts}
    \end{answer}
  \recommended \item 
    该向量在此集合中。
    参数取什么值会产生这个向量？
    \begin{exparts}
      \partsitem $\colvec[r]{5 \\ -5}$、
        $\set{\colvec[r]{1 \\ -1}k\suchthat k\in\Re}$
      \partsitem $\colvec[r]{-1 \\ 2 \\ 1}$、
        $\set{\colvec[r]{-2 \\ 1 \\ 0}i
           +\colvec[r]{3 \\ 0 \\ 1}j\suchthat i,j\in\Re}$
      \partsitem $\colvec[r]{0 \\ -4 \\ 2}$、
        $\set{\colvec[r]{1 \\ 1 \\ 0}m
               +\colvec[r]{2 \\ 0 \\ 1}n\suchthat m,n\in\Re}$
    \end{exparts}
    \begin{answer}
      对每一小题，我们通过考察分量的等式得到一个
      线性方程组。
      \begin{exparts}
       \partsitem $k=5$
       \partsitem 第二个分量表明$i=2$，第三个
       分量表明$j=1$。
       \partsitem $m=-4$，$n=2$
      \end{exparts} 
    \end{answer}
  \item 
    判断该向量是否在集合中。
    \begin{exparts}
      \partsitem $\colvec[r]{3 \\ -1}$、
        $\set{\colvec[r]{-6 \\ 2}k\suchthat k\in\Re}$
      \partsitem $\colvec[r]{5 \\ 4}$、
        $\set{\colvec[r]{5 \\ -4}j\suchthat j\in\Re}$
      \partsitem $\colvec[r]{2 \\ 1 \\ -1}$、
        $\set{\colvec[r]{0 \\ 3 \\ -7}
             +\colvec[r]{1 \\ -1 \\ 3}r\suchthat r\in\Re}$
      \partsitem $\colvec[r]{1 \\ 0 \\ 1}$、
        $\set{\colvec[r]{2 \\ 0 \\ 1}j
            +\colvec[r]{-3 \\ -1 \\ 1}k\suchthat j,k\in\Re}$
    \end{exparts}
    \begin{answer}
      对每一小题，我们通过考察分量的等式得到一个
      线性方程组。
      \begin{exparts}
        \partsitem 是；取$k=-1/2$。
        \partsitem 否；由方程$5=5\cdot j$和
            $4=-4\cdot j$组成的方程组无解。
        \partsitem 是；取$r=2$。
        \partsitem 否。
           第二个分量给出$k=0$。
           于是第三个分量给出$j=1$。
           但第一个分量对不上。
      \end{exparts}
     \end{answer}
  \item \cite{Cleary}
    一个有1200英亩土地的农民在考虑种植三种不同的庄稼：
    玉米、大豆和燕麦。   
    农民想用全部~$1200$英亩土地。  
    玉米种子每英亩花费\$$20$，而大豆种子和燕麦种子分别
    每英亩花费\$$50$和\$$12$。  
    农民有\$$40\,000$可用于购买种子，并打算
    把它全部花掉。
    \begin{exparts}  
      \item 利用上面的信息建立两个含三个未知数的线性方程
        并求解它。
     \item 该方程组的解就是农民可以做的选择。  
        写出两个合理的解。
     \item 假设秋天庄稼成熟时，农民
        每英亩玉米可带来\$$100$的收入，
        每英亩大豆\$$300$、每英亩燕麦\$$80$。  
        前一题的两个解中哪一个原本会带来
        更大的收入？ 
    \end{exparts}
    \begin{answer}
      \begin{exparts}
        \item 设$c$为玉米的英亩数，$s$为
          大豆的英亩数，$a$为燕麦的英亩数。
          \begin{equation*}
            \begin{linsys}{3}
              c   &+   &s   &+   &a   &=   &1200 \\ 
            20c   &+   &50s &+   &12a &=   &40\,000  
            \end{linsys}
            \grstep{-20\rho_1+\rho_2}
            \begin{linsys}{3}
              c   &+   &s   &+   &a   &=   &1200 \\ 
                  &    &30s &-   &8a  &=   &16\,000  
            \end{linsys}
          \end{equation*}
          为描述解集，我们可以用$a$参数化。
          \begin{equation*}
            \set{\colvec{c \\ s \\ a}
                 =\colvec{20\,000/30 \\ 16\,000/30 \\ 0}
                  +\colvec{-38/30 \\ 8/30 \\ 1}a
                 \suchthat a\in\Re}
          \end{equation*}
        \item 这里有许多可能的答案。 
          例如我们可以取$a=0$得到$c=20\,000/30\approx 666.66$和
          $s=16000/30\approx 533.33$。
          另一个例子是取$a=20\,000/38\approx 526.32$，给出
          $c=0$和$s=7360/38\approx 193.68$。
        \item 把前一题的答案代入
          $100c+300s+80a$。
      \end{exparts}
    \end{answer}
  \item 
    参数化这个单方程方程组的解集。
    \begin{equation*}
      x_1+x_2+\cdots+x_n=0
    \end{equation*}
    \begin{answer}
      这个方程组有一个方程。
      首变量是\( x_1 \)，其余变量是自由变量。
      \begin{equation*}
        \set{\colvec{-1 \\ 1 \\ \vdotswithin{-1} \\ 0}x_2
             +\cdots+
             \colvec{-1 \\ 0 \\ \vdotswithin{-1} \\ 1}x_n
             \suchthat x_2,\ldots,x_n\in\Re}
      \end{equation*}  
     \end{answer}
  \recommended \item 
    \begin{exparts}
    \partsitem 对左边应用高斯消元法来求解
      \begin{equation*}
        \begin{linsys}{4}
          x  &+  &2y  &    &    &-   &w   &=   &a   \\
         2x  &   &    &+   &z   &    &    &=   &b   \\
          x  &+  &y   &    &    &+   &2w  &=   &c   
        \end{linsys}
      \end{equation*}
      把\( x \)、\( y \)、\( z \)和\( w \)用
      常数$a$、$b$和$c$表示。
    \partsitem 用前一题的答案求解这个方程组。
      \begin{equation*}
        \begin{linsys}{4}
          x  &+  &2y  &    &    &-   &w   &=   &3   \\
         2x  &   &    &+   &z   &    &    &=   &1   \\
          x  &+  &y   &    &    &+   &2w  &=   &-2
        \end{linsys}
      \end{equation*}
    \end{exparts}
    \begin{answer}
      \begin{exparts}
        \partsitem 这里高斯消元法给出
          \begin{align*}
            \begin{amat}{4}
              1  &2  &0  &-1  &a  \\
              2  &0  &1  &0   &b  \\
              1  &1  &0  &2   &c
            \end{amat}
            &\grstep[-\rho_1+\rho_3]{-2\rho_1+\rho_2}
            \begin{amat}{4}
              1  &2  &0  &-1  &a  \\
              0  &-4 &1  &2   &-2a+b  \\
              0  &-1 &0  &3   &-a+c
            \end{amat}                                  \\
            &\grstep{-(1/4)\rho_2+\rho_3}
            \begin{amat}{4}
              1  &2  &0    &-1  &a  \\
              0  &-4 &1    &2   &-2a+b  \\
              0  &0  &-1/4 &5/2 &-(1/2)a-(1/4)b+c
            \end{amat}
          \end{align*}
          留下\( w \)为自由变量。
          求解：\( z=2a+b-4c+10w \)，
          且\( -4y=-2a+b-(2a+b-4c+10w)-2w \)于是
          \( y=a-c+3w \)，并且
          \( x=a-2(a-c+3w)+w=-a+2c-5w. \)
          因此解集如下。
          \begin{equation*}
             \set{\colvec{-a+2c \\ a-c \\ 2a+b-4c \\ 0}
                  +\colvec{-5 \\ 3 \\ 10 \\ 1}w
                  \suchthat w\in\Re}
          \end{equation*}
        \partsitem 代入\( a=3 \)、\( b=1 \)和\( c=-2 \)。
          \begin{equation*}
             \set{\colvec[r]{-7 \\ 5 \\ 15 \\ 0}
                  +\colvec[r]{-5 \\ 3 \\ 10 \\ 1}w
                  \suchthat w\in\Re}
          \end{equation*}
      \end{exparts}  
     \end{answer}
  \item 
    在矩阵元素的记法“\( a_{i,j} \)”中，
    为什么需要逗号？
    \begin{answer}
       省略逗号，比如写成\( a_{123} \)，
       是有歧义的，因为它可能指$a_{1,23}$或$a_{12,3}$。  
    \end{answer}
  \recommended \item 
    给出\( \nbyn{4} \)矩阵，其
    \( i,j \)元素为
    \begin{exparts*}
      \partsitem \( i+j \)；
      \partsitem \( -1 \)的\( i+j \)次幂。
    \end{exparts*}
    \begin{answer}
      \begin{exparts*}
        \partsitem \(
           \begin{mat}[r]
             2  &3  &4  &5  \\
             3  &4  &5  &6  \\
             4  &5  &6  &7  \\
             5  &6  &7  &8
           \end{mat} \)
        \partsitem \(
           \begin{mat}[r]
             1  &-1  &1   &-1  \\
            -1  &1   &-1  &1  \\
             1  &-1  &1   &-1  \\
            -1  &1   &-1  &1
           \end{mat} \)
      \end{exparts*}  
    \end{answer}
  \item  
    对任意矩阵\( A \)，
    \( A \)的\definend{转置}\index{transpose}
    \index{matrix!transpose}
    记作
    \( \trans{A} \)，是以\( A \)的行作为列的矩阵。
    求下面每一个的转置。
    \begin{exparts*}
      \partsitem \( \begin{mat}[r]
                  1  &2  &3  \\
                  4  &5  &6
               \end{mat}  \)
      \partsitem \( \begin{mat}[r]
                  2  &-3 \\
                  1  &1
               \end{mat}  \)
      \partsitem \( \begin{mat}[r]
                  5  &10 \\
                 10  &5
               \end{mat}  \)
      \partsitem \( \colvec[r]{1 \\ 1 \\ 0} \)
    \end{exparts*}
    \begin{answer}
      \begin{exparts*}
        \partsitem \( \begin{mat}[r]
                   1  &4  \\
                   2  &5  \\
                   3  &6
                 \end{mat}  \)
        \partsitem \( \begin{mat}[r]
                   2  &1  \\
                  -3  &1
                 \end{mat}  \)
        \partsitem \( \begin{mat}[r]
                   5  &10 \\
                  10  &5
                 \end{mat}  \)
        \partsitem \( \rowvec{1 &1 &0}  \)
      \end{exparts*}  
     \end{answer}
  \recommended \item 
    \begin{exparts}
      \partsitem 描述所有满足\( f(1)=2 \)且\( f(-1)=6 \)的
        函数\( f(x)=ax^2+bx+c \)。
      \partsitem 描述所有满足\( f(1)=2 \)的
        函数\( f(x)=ax^2+bx+c \)。
    \end{exparts}
    \begin{answer}
      \begin{exparts}
        \partsitem 代入\( x=1 \)和\( x=-1 \)得到
          \begin{equation*}
            \begin{linsys}{3}
              a  &+  &b   &+  &c  &=  &2  \\
              a  &-  &b   &+  &c  &=  &6  
            \end{linsys}
            \grstep{-\rho_1+\rho_2}
            \begin{linsys}{3}
              a  &+  &b   &+  &c  &=  &2  \\
                 &   &-2b &   &   &=  &4  
              \end{linsys}
          \end{equation*}
          所以函数的集合是
          \( \set{f(x)=(4-c)x^2-2x+c\suchthat c\in\Re} \)。
        \partsitem 代入\( x=1 \)得到
          \begin{equation*}
            \begin{linsys}{3}
              a  &+  &b   &+  &c  &=  &2  
            \end{linsys}
          \end{equation*}
          所以函数的集合是
          \( \set{f(x)=(2-b-c)x^2+bx+c\suchthat b,c\in\Re} \)。
      \end{exparts}  
    \end{answer}
  \item 证明：平面\( \Re^2 \)中任意五个点都位于一条共同的二次曲线上，也就是说，它们都满足某个形如
    \( ax^2+by^2+cxy+dx+ey+f=0 \)的方程，其中\( a,\,\ldots\,,f \)
    中有些不为零。
    \begin{answer}
      把五对$(x,y)$代入，我们得到一个有
      五个方程、六个未知数$a$、\ldots、$f$的方程组。
      由于未知数比方程多，如果方程之间没有不相容，
      那么就有无穷多解
      （至少有一个变量最终会成为自由变量）。

      但不相容不可能存在，因为$a=0$、\ldots、$f=0$是
      一个解（我们只是用这个零解来表明方程组
      是相容的\Dash 前面一段表明
      存在非零解）。 
    \end{answer}
  \item 
    编出一个四方程/四未知数的方程组，使其具有
    \begin{exparts}
      \partsitem 单参数的解集；
      \partsitem 双参数的解集；
      \partsitem 三参数的解集。
    \end{exparts}
    \begin{answer}
      \begin{exparts}
      \partsitem 这是一个\Dash 第四个方程是多余的
        但仍然可以。
        \begin{equation*}
          \begin{linsys}{4}
             x  &+  &y  &-  &z  &+  &w  &=  &0  \\
                &   &y  &-  &z  &   &   &=  &0  \\
                &   &   &   &2z &+  &2w &=  &0  \\
                &   &   &   &z  &+  &w  &=  &0
          \end{linsys}
        \end{equation*}
      \partsitem 这是一个。
        \begin{equation*}
          \begin{linsys}{4}
             x  &+  &y  &-  &z  &+  &w  &=  &0  \\
                &   &   &   &   &   &w  &=  &0  \\
                &   &   &   &   &   &w  &=  &0  \\
                &   &   &   &   &   &w  &=  &0
          \end{linsys}
        \end{equation*}
      \partsitem 这是一个。
        \begin{equation*}
          \begin{linsys}{4}
             x  &+  &y  &-  &z  &+  &w  &=  &0  \\
             x  &+  &y  &-  &z  &+  &w  &=  &0  \\
             x  &+  &y  &-  &z  &+  &w  &=  &0  \\
             x  &+  &y  &-  &z  &+  &w  &=  &0 
          \end{linsys}
        \end{equation*}
    \end{exparts}  
   \end{answer}
  \puzzle \item 
    \cite{Shepelev}
    这个谜题来自俄罗斯网站
    \texttt{http://www.arbuz.uz/}，它有很多种解法，
    但我的解法用到线性代数，而且非常朴素。有一颗行星上居住着arbuzoid
    （从俄语翻译过来即“西瓜人”）。
    这些生物有三种颜色：红、绿、蓝。
    有$13$~个红色的
    arbuzoid、$15$~个蓝色的、$17$~个绿色的。
    当
    两只颜色不同的arbuzoid相遇时，它们
    都变成第三种颜色。

    问题是，有没有可能出现它们
    全部变成同一种颜色的情况？
     \begin{answer}
        \answerasgiven
        我的解法是把arbuzoid的
        数量定义为$3$维向量，并把所有可能的
        基本转变也表示成这样的向量：
        \begin{center}
          \begin{tabular}{rr}
            R: &$13$  \\
            G: &$15$  \\
            B: &$17$
          \end{tabular}
          \qquad
          操作：
          $\colvec[r]{-1 \\ -1 \\ 2}$、
          $\colvec[r]{-1 \\ 2 \\ -1}$、
          以及
          $\colvec[r]{2 \\ -1 \\ -1}$
        \end{center}
        现在，只需检查下面某个线性方程组的解
        是否
        存在：
        \begin{equation*}
          \colvec[r]{13 \\ 15 \\ 17}
          +x\colvec[r]{-1 \\ -1  \\ 2}
          +y\colvec[r]{-1 \\ 2 \\ -1}
          +z\colvec[r]{2 \\ -1 \\ -1}
          =\colvec[r]{0 \\ 0 \\ 45}
          \qquad
          \text{（或 $\colvec[r]{0 \\ 45 \\ 0}$ 或 $\colvec[r]{45 \\ 0 \\ 0}$）}
        \end{equation*}
        求解
        \begin{equation*}
          \begin{amat}[r]{3}
           -1  &-1 &2  &-13  \\
           -1  &2  &-1 &-15  \\
            2  &-1 &-1 &28
          \end{amat}
          \grstep[2\rho_1+\rho_3]{-\rho_1+\rho_2}
          \repeatedgrstep{\rho_2+\rho_3}
          \begin{amat}[r]{3}
           -1  &-1 &2  &-13  \\
            0  &3  &-3 &-2  \\
            0  &0  &0  &0
          \end{amat}
        \end{equation*}
        给出$y+2/3=z$，所以如果变换次数$z$是一个整数
        那么$y$就不是。
        另外两个方程组给出类似的结论，因此无解。
     \end{answer}
  \puzzle \item 
    \cite{USSROlympiad174}
    \begin{exparts}
      \partsitem 求解方程组。
        \begin{equation*}
          \begin{linsys}{2}
            ax  &+  &y  &=  &a^2  \\
             x  &+  &ay &=  &1
         \end{linsys}
        \end{equation*}
        对哪些$a$值，方程组无解，而对哪些$a$值，方程组有无穷多解？
      \partsitem 对方程组回答上面的问题。
        \begin{equation*}
          \begin{linsys}{2}
            ax  &+  &y  &=  &a^3  \\    
             x  &+  &ay &=  &1
          \end{linsys}
        \end{equation*}
    \end{exparts}
    \begin{answer}
       \answerasgiven
       \begin{exparts}
        \partsitem 对方程组作形式求解得到
          \begin{equation*}
            x=\frac{a^3-1}{a^2-1}  
            \qquad
            y=\frac{-a^2+a}{a^2-1}.
          \end{equation*}
          如果$a+1\neq 0$且$a-1\neq 0$，那么方程组有唯一的
          解
          \begin{equation*}
            x=\frac{a^2+a+1}{a+1}
            \qquad
            y=\frac{-a}{a+1}.
          \end{equation*}
          如果$a=-1$，或者$a=+1$，那么这些公式没有意义；在
          第一种情形，我们得到方程组
          \begin{equation*}
            \left\{ 
            \begin{linsys}{2}
              -x &+  &y  &=  &1 \\
               x &-  &y  &=  &1
            \end{linsys}\right.
          \end{equation*}
          这是一个矛盾方程组。
          在第二种情形，我们有
          \begin{equation*}
            \left\{
            \begin{linsys}{2}
               x &+  &y  &=  &1 \\
               x &+  &y  &=  &1
            \end{linsys}\right.
          \end{equation*}
          它有无穷多个解（例如，$x$任取，
          $y=1-x$）。
        \partsitem 求解方程组得到
          \begin{equation*}
            x=\frac{a^4-1}{a^2-1}
            \qquad
            y=\frac{-a^3+a}{a^2-1}.
          \end{equation*}
          这里，若$a^2-1\neq 0$，则方程组有唯一解
          $x=a^2+1$，$y=-a$。
          对于$a=-1$和$a=1$，我们得到方程组
          \begin{equation*}
            \left\{
            \begin{linsys}{2}
              -x &+  &y  &=  &-1 \\
               x &-  &y  &=  &1
            \end{linsys}\right.
            \qquad
            \left\{
            \begin{linsys}{2}
               x &+  &y  &=  &1 \\
               x &+  &y  &=  &1
            \end{linsys}\right.
         \end{equation*}
         两者都有无穷多个解。
      \end{exparts}
    \end{answer}
  \puzzle \item 
    \cite{MathMag52p48}
    在空气中，一个表面镀金的球重\( 7588 \)
    克。
    已知它可能含有铝、铜、银或铅中的一种或多种金属。
    在标准条件下依次在水中、苯中、
    酒精中和甘油中称量时，其重量分别为\( 6588 \)、\( 6688 \)、
    \( 6778 \)和\( 6328 \)克。
    如果指定的各种物质的比重取如下数值，那么它含有上述金属各多少（如果有的话）？
    \begin{center}
       \begin{tabular}{lrclr}
         铝  &\( 2.7 \)
            &\makebox[3em]{\mbox{}\hfill\mbox{}} &酒精 &0.81 \\
         铜    &\( 8.9 \)  &     &苯   &\( 0.90 \) \\
         金      &\( 19.3 \) &     &甘油 &\( 1.26 \) \\
         铅      &\( 11.3 \) &     &水     &\( 1.00 \) \\
         银    &\( 10.8 \)
       \end{tabular}
    \end{center}
    \begin{answer}
      \answerasgiven
      设\( u \)、\( v \)、\( x \)、\( y \)、\( z \)分别为球中所含
      Al、Cu、Pb、Ag和Au的体积（单位
      \( {\rm cm}^3 \)），我们假设球不是空心的。
      由于在水中（比重\( 1.00 \)）的失重
      为
      \( 1000 \)克，球的体积为\( 1000\mbox{ cm}^3 \)。
      于是这些数据（其中一些是多余的，尽管是相容的）只导出
      \( 2 \)个独立的方程，一个联系体积，另一个联系重量。
      \begin{equation*}
        \begin{linsys}{5}
           u  &+  &v    &+  &x     &+  &y     &+  &z     &=  &1000  \\
        2.7u  &+  &8.9v &+  &11.3x &+  &10.5y &+  &19.3z &=  &7558
        \end{linsys}
      \end{equation*}
      显然，球必须含有一些铝，以使其平均比重
      低于其他所有金属的比重。
      问题的这一部分没有唯一的结果，因为三种金属的量
      可以任意选取，只要所选取的量
      不会使任何金属的量为负。

      如果球只含铝和金，那么有
      \( 294.5\mbox{ cm}^3 \)的金和\( 705.5\mbox{ cm}^3 \)的铝。
      另一种可能是Cu、Au、Pb和
      Ag各\( 124.7\mbox{ cm}^3 \)，以及\( 501.2\mbox{ cm}^3  \)的Al。   
    \end{answer}
\end{exercises}

























